\documentclass[11pt,a4paper]{article}

\usepackage[T1]{fontenc}
\usepackage[utf8]{inputenc}
\usepackage{lmodern}
\usepackage[margin=2.3cm]{geometry}
\usepackage{amsmath,amssymb}
\usepackage{booktabs}
\usepackage{array}
\usepackage{tabularx}
\usepackage{float}
\usepackage{caption}
\usepackage{tikz}
\usepackage{xcolor}
\usepackage{enumitem}
\usepackage{microtype}
\usepackage[numbers,sort&compress]{natbib}
\usepackage[hidelinks]{hyperref}

\captionsetup{font=small,labelfont=bf}
\setlist{nosep,leftmargin=1.5em}
\setlength{\parskip}{0.3em}

\newcommand{\maxseg}{\mathit{maxSeg}}
\newcommand{\MLU}{\mathrm{MLU}}

\title{Segment routing under planned link outages:\\
a MILP study of the Orange T-ASR problem}
\author{[Full name 1 -- Student ID] \and [Full name 2 -- Student ID]\\[0.3em]
\small Mathematical Modelling project, SoICT}
\date{October 2026}

\begin{document}
\maketitle

\begin{abstract}
\noindent
We study the T-Adaptive Segment Routing problem proposed by Orange for the EURO/ROADEF 2026
challenge. Each traffic demand of an IP network has to be assigned a segment path, i.e.\ a short
list of waypoints between which packets follow shortest paths, so that the most loaded link stays
as low as possible, both in the nominal state and when a link is switched off for maintenance.
Following the course guidelines we build a mixed-integer linear model for the nominal period, then
extend it to two periods with a budget on the number of routing changes. Split ratios are computed
from shortest-path counts. The model is solved with JuMP and HiGHS on the 20 instances of set~A
under a 15\,GiB memory limit. The full model can only be built for the five smallest instances, so
for the others we solve a restricted model with a small set of candidate waypoints per demand, and
we use a multi-commodity flow relaxation as an independent lower bound. Segment routing lowers the
maximum link utilization by about 21\,\% on average and by up to 80\,\%. Sixteen of the 40 runs are
proven optimal for the original problem.
\end{abstract}

%======================================================================
\section{Introduction}
%======================================================================

Orange operates large IP/MPLS backbones with thousands of routers and tens of thousands of traffic
demands. In such networks packets are usually forwarded along shortest paths computed by the
interior gateway protocol from administrative link weights. When several shortest paths exist, the
traffic is split between them (Equal-Cost Multi-Path, ECMP). Routing on shortest paths is simple to
operate, but the only way to steer traffic is to change the weights, and finding weights that avoid
congestion is a hard problem in its own right~\cite{fortz2000}. A weight change also moves many
demands at once, which is not what an operator wants during a maintenance window.

Segment Routing (SR)~\cite{rfc8402} gives a finer control. The ingress router writes into the
packet header a list of intermediate nodes, called waypoints or node segments. The packet travels
on shortest paths from the source to the first waypoint, then to the next one, and so on until the
destination. With one or two waypoints a demand can be pushed away from a congested area without
touching the link weights~\cite{bhatia2015}. The price is that every waypoint is a configuration
item that someone has to deploy and maintain.

The challenge problem, called T-Adaptive Segment Routing (T-ASR)~\cite{orange2026}, adds a time
dimension. Routing is planned over a horizon of periods. Period~0 is the nominal network; in later
periods some links are taken down for planned interventions, and the traffic matrix may change as
well. For each period one has to choose a segment path for every demand, with at most $\maxseg$
segments, while the number of segment changes between consecutive periods is bounded by a budget.
The official objective minimizes all link loads in lexicographic order.

This report follows the steps given for the course project. We first compute the split ratios from
the network data, then write and solve a model for period~0 alone (no budget), and finally extend
it to the two periods of the set~A instances with the budget constraint. Section~\ref{sec:problem}
describes the data, Section~\ref{sec:ratios} the split ratios and Section~\ref{sec:model} the model.
Section~\ref{sec:method} explains how we solved it on large instances with limited memory, and
Section~\ref{sec:exp} reports the experiments. We close with remarks on scalability and on what we
would do next.

%======================================================================
\section{Problem description and data}\label{sec:problem}
%======================================================================

The network is a directed graph $G=(V,A)$ with $n=|V|$ routers. Each arc $a\in A$ has a weight
$\omega(a)>0$, used to compute shortest paths, and a capacity $c(a)>0$. A demand $d\in D$ is a pair
$(s_d,t_d)$ with a traffic volume $\nu(d,t)\ge 0$ in each period $t\in T=\{0,\dots,h-1\}$. An
intervention scenario gives, for each $t\ge 1$, the set $q(t)\subseteq A$ of arcs that are down, so
the network of period $t$ is $G_t=(V,A\setminus q(t))$. A segment path of $d$ is a sequence
$\langle s_d,w_1,\dots,w_\ell,t_d\rangle$; its segments are the consecutive pairs $(s_d,w_1),
\dots,(w_\ell,t_d)$, and their number $\ell+1$ may not exceed $\maxseg$. The distance between two
routing schemes is the number of segments that appear in exactly one of them, and it may not
exceed the budget $\kappa(t)$ between periods $t-1$ and $t$.

Table~\ref{tab:data} lists where each parameter comes from. The set~A benchmark has 20 instances
with 20 to 400 nodes, 80 to 2000 arcs and 20 to 6000 demands (Table~\ref{tab:instances}). All of
them have two periods, $\maxseg=6$, exactly one link down in period~1, and a budget between 1 and
90. A few details of the files mattered for the implementation. Nine instances carry the flag
\texttt{multigraph: true}, but none of them actually contains parallel arcs, so a simple directed
graph is enough. In setA-08 and setA-14 the metrics are not integers, which affects the counting of
shortest paths (Section~\ref{sec:ratios}). Finally, the volumes of period~1 are not a scaled copy of
period~0; some demands grow by a factor of more than 60, so period~1 differs from period~0 by more than the
missing link.

\begin{table}[tbp]
\centering\small
\caption{Model parameters and their source files.}\label{tab:data}
\begin{tabular}{@{}lll@{}}
\toprule
Symbol & Meaning & Source \\
\midrule
$c(a)$, $\omega(a)$ & capacity and IGP metric of arc $a$ & \texttt{net.json}, fields \texttt{capacity}, \texttt{metric}\\
$r(i,j,a,t)$ & share of the $i\to j$ flow carried by $a$ in $G_t$ & computed from the metrics (Sec.~\ref{sec:ratios})\\
$\nu(d,t)$ & volume of demand $d$ in period $t$ & \texttt{tm.json}, field \texttt{v[t]}\\
$\maxseg$ & maximum number of segments & \texttt{scenario.json}, \texttt{max\_segments}\\
$q(t)$, $\kappa(t)$ & arcs down and change budget in period $t$ & \texttt{scenario.json}, \texttt{interventions}, \texttt{budget}\\
\bottomrule
\end{tabular}
\end{table}

%======================================================================
\section{Split ratios}\label{sec:ratios}
%======================================================================

The split ratio $r(i,j,a,t)$ is the fraction of the traffic sent from $i$ to $j$ that crosses arc
$a$ in $G_t$. It is defined for every ordered pair $(i,j)$, since any pair may become a segment.
Following the course notes we take
\begin{equation}
  r(i,j,a,t)=\frac{\sigma^a_{ij}(t)}{\sigma_{ij}(t)},
\end{equation}
where $\sigma_{ij}$ is the number of shortest $i\to j$ paths and $\sigma^a_{ij}$ the number of
those that use $a$. Both are obtained without enumerating paths. A Dijkstra run from $i$ keeps,
next to the distance label $d(v)$, a counter $\sigma(v)$: when an arc $(u,v)$ gives a strictly
shorter path, $\sigma(v)$ is reset to $\sigma(u)$, and when it gives a path of equal length,
$\sigma(u)$ is added to $\sigma(v)$. Running it from every node yields the matrices $d_{ij}$ and
$\sigma_{ij}$; in Julia this is \texttt{dijkstra\_shortest\_paths} with \texttt{allpaths=true}.
An arc $a=(u,v)$ lies on a shortest $i\to j$ path if and only if $d_{iu}+\omega(a)+d_{vj}=d_{ij}$.
Such a path is a shortest $i\to u$ path, followed by $a$, followed by a shortest $v\to j$ path, and
every combination of the two parts is a shortest path, so $\sigma^a_{ij}=\sigma_{iu}\sigma_{vj}$.
Hence
\begin{equation}
  r(i,j,a,t)=
  \begin{cases}
    \sigma_{iu}\,\sigma_{vj}/\sigma_{ij} & \text{if } d_{iu}+\omega(a)+d_{vj}=d_{ij},\\
    0 & \text{otherwise.}
  \end{cases}
\end{equation}
The arcs with $r>0$ form the forwarding graph $FG(i,j)$. For every pair, $r(i,j,\cdot,t)$ must be a
unit flow from $i$ to $j$, and we check this flow conservation on every pair of every instance.
The check caught a numerical problem. With non-integer metrics, two paths of the same length can
differ in the last bit of their floating point sum, and the exact comparison used by Graphs.jl then
misses some of them. On setA-08 and setA-14 the conservation error reached 1, meaning that part of
the flow simply disappeared. In that case we recompute the counts with our own Dijkstra in which
two lengths are equal when they agree up to a relative tolerance of $10^{-9}$. After this the error
is zero on all instances. The ratios are stored sparsely, pair by pair; for setA-20 there are
716\,276 non-zero values in period~0 and the whole computation takes less than 4 seconds.

This definition differs from the one in the challenge description, where ECMP splits the flow
evenly at every node of the forwarding graph. On the toy instance of the subject, for the demand
$(v_0,v_5)$ there are three shortest paths, one through $v_3$ only and two through $v_4$. Path
counting gives $r=1/3$ to arc $(v_1,v_3)$ and $2/3$ to $(v_1,v_4)$, while per-node ECMP gives $1/2$
to each. For the demand $(v_2,v_5)$ both rules agree. We use path counting as required by the
course; with it the toy instance has a no-waypoint MLU of $0.4167$ instead of the $0.4375$ quoted in
the subject. To make sure the rest of the code is right, our test suite also implements the per-node
rule and checks that the load computation then reproduces the figures of the subject exactly.

%======================================================================
\section{Mathematical model}\label{sec:model}
%======================================================================

Let $P_t=\{(i,j)\in V^2: i\neq j,\ \sigma_{ij}(t)>0\}$ be the pairs that can be used as segments in
period $t$. The decision variables are
\begin{itemize}
  \item $x^{dt}_{ij}\in\{0,1\}$, equal to 1 if the segment path of demand $d$ in period $t$ contains
  the segment $(i,j)$, for $(i,j)\in P_t$;
  \item $U\ge 0$, the maximum link utilization over all arcs and periods.
\end{itemize}
The model reads
\begin{align}
\min\quad & U \nonumber\\
\text{s.t.}\quad
& \sum_{j:(i,j)\in P_t} x^{dt}_{ij}-\sum_{j:(j,i)\in P_t} x^{dt}_{ji}=
  \begin{cases}1 & i=s_d\\ -1 & i=t_d\\ 0 & \text{otherwise}\end{cases}
  && \forall d\in D,\ i\in V,\ t\in T, \label{c:flow}\\
& \sum_{(i,j)\in P_t} x^{dt}_{ij}\le \maxseg && \forall d\in D,\ t\in T, \label{c:seg}\\
& \sum_{d\in D}\sum_{(i,j)\in P_t} r(i,j,a,t)\,\nu(d,t)\,x^{dt}_{ij}\le c(a)\,U
  && \forall t\in T,\ a\in A\setminus q(t), \label{c:load}\\
& \sum_{d\in D}\sum_{(i,j)} \bigl|x^{dt}_{ij}-x^{d,t-1}_{ij}\bigr|\le \kappa(t)
  && \forall t\in T\setminus\{0\}, \label{c:budget}\\
& x^{dt}_{ij}\in\{0,1\},\quad U\ge 0. \nonumber
\end{align}
For step~6 of the project $T=\{0\}$ and constraint~\eqref{c:budget} is absent; for step~9
$T=\{0,1\}$.

\paragraph{Why each constraint.}
Think of every pair $(i,j)\in P_t$ as a virtual arc of a complete graph. Constraint~\eqref{c:flow}
asks the virtual arcs chosen for $d$ to form a path from $s_d$ to $t_d$; the inner nodes of that path
are the waypoints. The number of virtual arcs on the path is the number of segments, which is what
\eqref{c:seg} limits. In~\eqref{c:load} the traffic of $d$ on a segment $(i,j)$ is spread over the
real arcs of $FG(i,j)$ according to $r$, and the sum over demands and segments is the flow on arc
$a$. Dividing by $c(a)$ gives the load $\lambda(a,t)$, so~\eqref{c:load} says $U\ge\lambda(a,t)$ for
every arc that is up; at the optimum $U$ equals the maximum load. Arcs in $q(t)$ get no constraint,
and they never carry flow anyway because they are absent from $G_t$. Constraint~\eqref{c:budget} is
the distance between the two routing schemes as defined in the subject.

\paragraph{Objective.}
The subject asks for the lexicographic minimum of all loads. Its first component is the maximum
load, so minimizing $U$ is the first level of that objective. Computing the full lexicographic
optimum would require a sequence of up to $|A|\,|T|$ MILPs, each as large as ours, which is out of
reach here. MLU is also the criterion most used in traffic engineering. With two periods the same
variable $U$ bounds the loads of both, so the objective is $\max(\MLU_0,\MLU_1)$.

\paragraph{Restricting the pairs.}
We only create $x^{dt}_{ij}$ for $(i,j)\in P_t$. If $j$ cannot be reached from $i$, all ratios of
$(i,j)$ are zero, and a variable on that pair would carry traffic for free. This happens in period~1
when the down arc disconnects some pairs, so the variables of such a pair are fixed to 0 in that
period.

\paragraph{Cycles.}
Constraint~\eqref{c:flow} allows, besides the $s_d$--$t_d$ path, cycles of virtual arcs that are
disconnected from it. They are harmless. Since $r\ge 0$ and $\nu\ge 0$, removing a cycle can only
decrease loads and the number of segments, so an optimal solution without cycles always exists and
no subtour constraint is needed. When reading a solution we follow the chosen segments from $s_d$
and drop anything not on the path.

\paragraph{Linearization of the budget.}
The usual way is to add $z\ge x^{dt}-x^{d,t-1}$ and $z\ge x^{d,t-1}-x^{dt}$ for each term and bound
$\sum z$. We use the identity $|\Delta|=2\Delta^+-\Delta$ instead, which needs one row per term:
\begin{equation}
  z^d_{ij}\ge x^{d1}_{ij}-x^{d0}_{ij},\quad z^d_{ij}\ge 0,\quad
  b_d=\sum_{(i,j)}\bigl(2z^d_{ij}-x^{d1}_{ij}+x^{d0}_{ij}\bigr),\quad
  \sum_{d\in D} b_d\le\kappa(1).
\end{equation}
Since $z$ is only bounded from below, a feasible $z$ exists exactly when $\sum|\Delta|\le\kappa$, so
the formulation is exact. The per-demand counters $b_d$ look redundant but they matter in practice.
Written as a single row, the budget constraint contains every $z$ variable of the model (hundreds of
thousands of coefficients), and the HiGHS presolve got stuck on it: on setA-20 it ran for more than
20 minutes without honouring a 60-second time limit. With the $b_d$ split, presolve takes about
5 seconds, and the step~9 model of setA-13, which had also blocked in presolve, is solved to
optimality in 38 seconds.

\paragraph{Aggregated load rows.}
Written as in~\eqref{c:load}, the load rows have $|D|\cdot\sum_{(i,j)}|FG(i,j)|$ non-zeros, for
example 23 million for setA-08. We introduce one continuous variable per segment and period,
$y^t_{ij}=\sum_d \nu(d,t)\,x^{dt}_{ij}$, and write the load rows as
$\sum_{(i,j)} r(i,j,a,t)\,y^t_{ij}\le c(a)U$. This is only a renaming, but the count drops to
$|D|\,|P_t|+\sum|FG(i,j)|$, i.e.\ 4.6 million for setA-08. Demands with zero volume in every period
are removed from the model and routed directly.

\paragraph{Size.}
The model has $|D|\cdot|P_t|\le |D|\,n(n-1)$ binary variables per period. The last column of
Table~\ref{tab:instances} gives this number for step~9: from 30\,400 for setA-01 to 1.9 billion for
setA-20. Size, not difficulty, is the first obstacle on this benchmark.

%======================================================================
\section{Solution approach}\label{sec:method}
%======================================================================

All experiments had to fit in 15\,GiB of memory. Measured on our runs, the full model takes roughly
2.5\,KB per variable once HiGHS has loaded it, so only instances with at most a few million
variables can be built. We used three ingredients.

\paragraph{Full model when it fits.}
When the full model has at most 2.5 million $x$ variables (setA-01 to setA-05, except setA-04 in
step~9, see below), we solve it directly. Two changes to the default settings were needed. First,
the LP relaxation of these models is large and highly degenerate. On setA-03 step~9 the dual simplex
performed 752\,481 iterations in one hour without finishing the root LP, whereas the interior point
method solved it in 36 seconds; we therefore set HiGHS to use the interior point method for the LP
at the root. Second, the primal heuristics of HiGHS rarely found anything better than the starting
solution on these models. We start from the solution of the restricted model described next, solved
with at most 10\,\% of the time limit. It is feasible for the full model and usually already very
good.

\paragraph{Restricted model otherwise.}
For larger instances, demand $d$ may only use waypoints from a candidate set $W_d$ of $K$ nodes,
those with the smallest detour $d_{s_dw}+d_{wt_d}-d_{s_dt_d}$. The segments of $d$ are the pairs of
$\{s_d\}\cup W_d\cup\{t_d\}$, without arcs entering $s_d$ or leaving $t_d$. We choose
$K=\lfloor\sqrt{600\,000/(2|D|)}\rfloor-1$, clipped to $[3,30]$, which gives about 600\,000 binary
variables in step~9; the same $K$ is used in step~6, so that the step~6 model is exactly the
period-0 part of the step~9 model. Every solution of the restricted model is a valid routing, so its
objective is an upper bound for the original problem. Its ``optimal'' status, however, only refers
to the restricted space, and its dual bound says nothing about the original problem.

\paragraph{An independent lower bound.}
To judge the restricted solutions we solve a multi-commodity flow relaxation: each demand may split
its traffic arbitrarily over all paths of $G_t$, and we minimize the maximum load. Any segment
routing is such a flow, so the optimum of this LP (aggregated by source node, solved by the interior
point method) is a valid lower bound on the T-ASR optimum. When a solution reaches this bound it is
optimal for the original problem. The bound can be weak, because SR cannot split a demand freely.

\paragraph{Correctness checks.}
Every solution is decoded into waypoint lists and checked by code that does not use the model: each
path is simple, has at most $\maxseg$ segments, uses only pairs reachable in $G_t$, the budget is
respected, and the MLU recomputed from the waypoints matches the objective value. The whole pipeline
was first tested on the toy instance of the subject, where all values can be checked by hand.

%======================================================================
\section{Numerical experiments}\label{sec:exp}
%======================================================================

\subsection{Setting}

The code is written in Julia 1.13.2 with JuMP 1.32.1~\cite{jump2023}, HiGHS 1.15.1 (through
HiGHS.jl 1.26.0)~\cite{highs2018}, Graphs.jl 1.15.0 and JSON3. Runs were made on a Linux server with
an AMD EPYC 7742 processor. HiGHS used a single thread, random seed 1 and a relative gap tolerance of
$10^{-4}$. The time limit was 1800\,s for step~6 and 3600\,s for step~9. Each solve ran in its own
process, and all processes together were confined to a Linux control group with a hard limit of
15\,GiB, several solves sharing it at the same time. CPU time is measured around the call to
\texttt{optimize!} only, after a warm-up run that removes Julia compilation time; it excludes the
time to compute the ratios (at most 4\,s) and to build the model (at most 95\,s).

The project guidelines ask for averages over ten runs. Because of the available time we made a single
run per instance and step, so the tables report one value and no standard deviation. Most instances
are solved at the root node, where the seed has little influence, but the runs that hit the time
limit could vary from one seed to another.

Table~\ref{tab:instances} describes the instances. Two observations help reading the results. The
no-waypoint MLU (``baseline'') is 1.0000 or very close to it for most instances and periods, which
suggests that the volumes were scaled so that shortest-path routing just saturates the network; any
value below 1 means that SR has removed the congestion. The multi-commodity flow bound is sometimes
very low (0.04 for setA-13, 16 and 19), which already tells us that it will not certify much on
those instances.

\begin{table}[tbp]
\centering\small
\caption{Set~A instances. $\kappa$: budget of period 1; down: id of the link removed in period 1;
baseline: MLU without waypoints; MCF: multi-commodity flow lower bound (step~6: period 0; step~9:
both periods); $|x|$: binary variables of the full step~9 model, in millions.}\label{tab:instances}
\setlength{\tabcolsep}{5pt}
\begin{tabular}{@{}rrrrrrcccccr@{}}
\toprule
 & & & & & & \multicolumn{2}{c}{Baseline} & \multicolumn{2}{c}{MCF bound} & \\
\cmidrule(lr){7-8}\cmidrule(lr){9-10}
Inst & $|V|$ & $|A|$ & $|D|$ & $\kappa$ & down & $t=0$ & $t=1$ & step 6 & step 9 & $|x|$ (M)\\
\midrule
\csname @@input\endcsname tables/instances.tex 
\bottomrule
\end{tabular}
\end{table}

\subsection{Period 0 only (step 6)}

Table~\ref{tab:step6} gives the results for period~0. Status, gap and node count are those reported
by HiGHS for the model that was solved, with gap $=1-\text{LB}/\text{UB}$. A star marks the solutions
that are provably optimal for the original problem, either because the full model was solved to
optimality or because the objective equals the MCF bound.

\begin{table}[tbp]
\centering\small
\caption{Step 6, period 0 only, time limit 1800\,s. Model: full model (exact) or restricted model
with $K$ candidate waypoints. $^\star$ provably optimal for the original problem. $^\ddagger$ HiGHS
exceeded the time limit.}\label{tab:step6}
\setlength{\tabcolsep}{4.5pt}
\begin{tabular}{@{}rrrrlllrrrcc@{}}
\toprule
Inst & $|V|$ & $|A|$ & $|D|$ & Model & Status & Best obj. & Gap (\%) & CPU (s) & Nodes & Baseline & MCF \\
\midrule
\csname @@input\endcsname tables/step6.tex 
\bottomrule
\end{tabular}
\end{table}

The full model is solved to optimality on the three smallest instances, in less than 7 minutes. On
setA-04 and setA-05 (475\,000 and 990\,000 binaries) HiGHS spends the whole time limit on the root
node. For setA-05 the starting solution from the restricted model already brings the MLU from 0.333
to 0.132, and the root bound shows that it is within 12.8\,\% of the optimum. Note that HiGHS overran
the limit on this instance by about 20 minutes; the interior point method only checks the clock
between iterations, and one iteration on a model of this size takes long. On setA-04 the start
phase found nothing better than the baseline within its 180 seconds, and neither did the full
model.

The restricted model behaves better than its size would suggest. Thirteen of the fifteen restricted
runs end with status optimal, all but one at the root node: on this problem the vertex returned by the
simplex method is very often already integral. Over the 20 instances the MLU decreases by 21.7\,\% on
average, by more than 50\,\% on setA-05, 10, 14 and 17. Seven results are provably optimal for the
original problem.

\subsection{Two periods with the budget (step 9)}

\begin{table}[tbp]
\centering\small
\caption{Step 9, periods 0 and 1 with the budget constraint, time limit 3600\,s. Last column:
number of changed segments over the budget $\kappa$. $^\star$ provably optimal for the original
problem (full model optimal, objective equal to the MCF bound, or equal to the proven step~6
optimum, which is a lower bound since step 6 is a relaxation of step 9).}\label{tab:step9}
\setlength{\tabcolsep}{3.6pt}
\begin{tabular}{@{}rrrrlllrrrccr@{}}
\toprule
Inst & $|V|$ & $|A|$ & $|D|$ & Model & Status & Best obj. & Gap (\%) & CPU (s) & Nodes & Baseline & MCF & Chg. \\
\midrule
\csname @@input\endcsname tables/step9.tex 
\bottomrule
\end{tabular}
\end{table}

Table~\ref{tab:step9} and Figure~\ref{fig:bars} show the two-period results. Removing period~1 and
the budget from the step~9 model leaves exactly the step~6 model, so a proven step~6 optimum is a
lower bound for step~9. This certifies setA-03: HiGHS reports a 33\,\% gap, but the solution reaches
the step~6 optimum 0.9435 and is therefore optimal. The same argument does not apply to restricted
models, whose step~6 value is only an upper bound. Altogether nine step~9 solutions are provably
optimal. The full model of setA-04 needed more than 9\,GB after half an hour and was stopped by our
memory guard; it was solved again with the restricted model ($K=20$).

The step~9 objective is never below the step~6 one, as expected, and the maximum is attained in
period~1 on most instances because of the missing link and the different traffic matrix. The budget
is rarely tight. In 11 instances the routing does not change at all between the two periods; only
setA-13 comes close to its budget (9 changes out of 12). On setA-05, with $\kappa=1$, the solution
uses no change at all and still lowers the MLU from 1 to 0.205, which matches the MCF bound.

\begin{figure}[tbp]
\centering
\def\barsdata{0/01/1.0000/0.9294,1/02/1.0981/0.9031,2/03/1.0000/0.9435,3/04/1.0000/0.7421,4/05/1.0000/0.2050,5/06/1.0000/1.0000,6/07/1.0000/0.9080,7/08/1.0000/1.0000,8/09/1.0000/0.8497,9/10/1.0000/0.4239,10/11/1.0000/0.7858,11/12/1.0001/1.0001,12/13/1.0000/0.6496,13/14/1.0000/0.5331,14/15/1.0000/0.8987,15/16/1.0000/0.8811,16/17/1.0001/0.4242,17/18/1.0000/1.0000,18/19/1.0000/0.7377,19/20/1.0000/0.9913}

\begin{tikzpicture}[x=0.72cm,y=3.2cm]
  \draw[gray!40] (-0.6,0) -- (19.6,0);
  \foreach \yy in {0.25,0.5,0.75,1} {
    \draw[gray!25] (-0.6,\yy) -- (19.6,\yy);
    \node[left,font=\scriptsize] at (-0.6,\yy) {\yy};
  }
  \foreach \k/\lab/\base/\obj in \barsdata {
    \fill[gray!35] (\k-0.38,0) rectangle (\k,\base);
    \fill[blue!65!black] (\k,0) rectangle (\k+0.38,\obj);
    \node[below,font=\scriptsize] at (\k,0) {\lab};
  }
  \fill[gray!35] (13.2,1.17) rectangle (13.5,1.24); \node[right,font=\scriptsize] at (13.5,1.205) {no waypoint};
  \fill[blue!65!black] (16.4,1.17) rectangle (16.7,1.24); \node[right,font=\scriptsize] at (16.7,1.205) {segment routing};
  \node[rotate=90,font=\scriptsize] at (-2.1,0.55) {MLU};
\end{tikzpicture}
\caption{Step 9: maximum link utilization over both periods without waypoints and with the best
segment routing found, per instance.}\label{fig:bars}
\end{figure}

\subsection{Why some instances do not improve}

Ten of the 40 runs end at the baseline value. They fall into three groups, and telling them apart
was one of the most useful parts of the analysis.

\emph{The baseline is already optimal.} On setA-07 (step~6) and setA-18 (both steps) the baseline
equals the MCF bound, so no routing at all can do better. In setA-18 the most loaded arc leaves
node 138, which has a single outgoing link; every demand that starts there has to use it, whatever
the waypoints.

\emph{The candidate set misses the useful waypoint.} On setA-12 the restricted model is solved to
optimality with MLU 1.0001 while the MCF bound is 0.58. The saturated arc is one of the two arcs
entering node 175. The other one starts at node 161 and has a metric about 27 times larger than the
surrounding arcs, so node 161 has a large detour and ranks 161st out of 198 candidates for the
demands involved. Routing just two of these demands through node 161 by hand gives an MLU of
0.9898. The restricted model is therefore optimal in its own space but misses a better solution; a
candidate rule that looks at link loads and not only at detours would fix this case.

\emph{The solver does not finish the root LP.} This concerns setA-04 in step~6 (see above) and
setA-06 and setA-08 in both steps ($K=23$ and $K=30$). On the last two the dual simplex was still at the root node when the time limit expired, with no useful
bound, so the starting solution was returned. We reran these four cases with the interior point
method at the root (Table~\ref{tab:ipm}). It always produced a bound, and on setA-08 step~6 it found
an optimal solution with MLU 0.3189 in under three minutes, a 68\,\% reduction. The interior point
method is not better everywhere, though: on setA-20 step~9 it reached the bound 0.9913 quickly but
found no integer solution below 1 in 20 minutes, while the simplex method solved the same model in
62 seconds. A simplex vertex is often integral on this problem, an interior point is not. Running
the simplex first and switching to the interior point method when the root LP stalls would combine
the two.

\begin{table}[tbp]
\centering\small
\caption{Supplementary runs on the two instances where the dual simplex did not finish the root LP:
same restricted model, root LP solved by dual simplex (main tables) or by the interior point method.}
\label{tab:ipm}
\begin{tabular}{@{}rrlccrccr@{}}
\toprule
 & & & \multicolumn{3}{c}{Dual simplex} & \multicolumn{3}{c}{Interior point}\\
\cmidrule(lr){4-6}\cmidrule(lr){7-9}
Inst & Step & Model & Best obj. & Bound & CPU (s) & Best obj. & Bound & CPU (s)\\
\midrule
\csname @@input\endcsname tables/ipm.tex 
\bottomrule
\end{tabular}
\end{table}

%======================================================================
\section{Scalability and practical use}
%======================================================================

The number of binary variables grows as $|D|\,n^2\,|T|$, and memory, not time, is the first limit.
With 15\,GiB the full model is usable up to about 100 nodes and 100 to 200 demands, which is the
size of a regional or metropolitan network, or of a reduced backbone used for planning studies. In
that range it gives proven optimal routings, or at least a good solution with a bound, within the
time limits we used. The subject allows routing plans to be computed offline, days or weeks ahead,
so a run of half an hour to an hour is acceptable in practice.

For national backbones the full model is out of reach, and the restricted model is a reasonable
compromise: it produced good solutions for all instances up to 400 nodes and 6000 demands, in a few
minutes and with less than 8\,GB. Its weak point is the choice of candidates, as setA-12 shows. Two
directions seem natural. The first is a column generation scheme in which segment paths are
generated by pricing~\cite{jadin2019}; this keeps the model small without fixing the candidates in
advance. The second, cheaper one is to grow the candidate set iteratively around the most loaded
arcs. Operationally, our results also suggest that few waypoints are needed: in the solutions most
demands use zero, one or two waypoints, and between the two periods the number of changes is far
below the budget.

%======================================================================
\section{Conclusion}
%======================================================================

We modelled the T-ASR problem as a MILP whose binary variables select segments for each demand and
period, with a flow formulation on the complete graph of segments, a segment count limit, load
constraints built from shortest-path split ratios, and a linearized change budget. Some reformulations
that do not change the model had a large effect on solvability: aggregating the load rows by segment,
and splitting the budget row by demand. On the 20 instances of set~A, segment routing reduced the
maximum link utilization by about 21\,\% on average and by up to 80\,\%. Sixteen of the 40 runs are
proven optimal; the others are feasible routings with either a solver gap or, for the restricted
model, a comparison with the multi-commodity flow bound.

The main limitations are the size of the full model, the heuristic choice of waypoint candidates for
large instances, the use of the maximum load instead of the full lexicographic objective, and the
single run per instance. The analysis of the failed cases points to concrete improvements: a
load-aware candidate rule, a fallback to the interior point method when the root LP stalls, and warm
starting the two-period model from the one-period solution.

\begin{thebibliography}{9}
\setlength{\itemsep}{0pt}
\small

\bibitem{orange2026}
A.~Benhamiche, Y.~Carlinet, M.~Chopin, E.~Gourdin, N.~Perrot.
\newblock Challenge EURO/ROADEF 2026, ``Keep the Flow!'': problem description.
\newblock Orange Research, 2025.

\bibitem{rfc8402}
C.~Filsfils, S.~Previdi, L.~Ginsberg, B.~Decraene, S.~Litkowski, R.~Shakir.
\newblock Segment Routing Architecture.
\newblock RFC 8402, IETF, 2018.

\bibitem{fortz2000}
B.~Fortz, M.~Thorup.
\newblock Internet traffic engineering by optimizing OSPF weights.
\newblock In \emph{Proc.\ IEEE INFOCOM}, pp.~519--528, 2000.

\bibitem{bhatia2015}
R.~Bhatia, F.~Hao, M.~Kodialam, T.~V.~Lakshman.
\newblock Optimized network traffic engineering using segment routing.
\newblock In \emph{Proc.\ IEEE INFOCOM}, pp.~657--665, 2015.

\bibitem{jadin2019}
M.~Jadin, F.~Aubry, P.~Schaus, O.~Bonaventure.
\newblock CG4SR: Near optimal traffic engineering for segment routing with column generation.
\newblock In \emph{Proc.\ IEEE INFOCOM}, pp.~1333--1341, 2019.

\bibitem{highs2018}
Q.~Huangfu, J.~A.~J.~Hall.
\newblock Parallelizing the dual revised simplex method.
\newblock \emph{Mathematical Programming Computation}, 10(1):119--142, 2018.

\bibitem{jump2023}
M.~Lubin, O.~Dowson, J.~Dias~Garcia, J.~Huchette, B.~Legat, J.~P.~Vielma.
\newblock JuMP 1.0: recent improvements to a modeling language for mathematical optimization.
\newblock \emph{Mathematical Programming Computation}, 15:581--589, 2023.

\end{thebibliography}

\end{document}
